% Vietnamese translation for OI Public Library
% Source-File: IOI中国国家候选队论文/2008/Day2/10.周冬《两极相通——浅析最大—最小定理在信息学竞赛中的应用》/冬令营论文演示文稿.ppt
\documentclass[11pt]{article}
\usepackage{oipl-vn}

\newcommand{\Oh}{\mathcal{O}}

\title{Bản trình chiếu: Định lý cực đại--cực tiểu trong thi tin học}
\author{Zhou Dong}
\date{2008}

\begin{document}
\maketitle

\origtitle{Slide companion for maximum--minimum theorems in informatics contests}
\sourcepath{IOI China candidate team papers / 2008 / Day 2 / Zhou Dong.ppt}

\section{Phạm vi bản dịch}

Tệp gốc chỉ có PowerPoint và phần chữ trích xuất được không đầy đủ. Các dấu
hiệu đọc được gồm những ví dụ như \textit{Muddy Fields}, \textit{Moving the
Hay}, thuật ngữ từ lý thuyết đồ thị, luồng mạng, tổ hợp, Dijkstra và các công
thức đối ngẫu. Bản ghi chú này ghi lại khung chủ đề chắc chắn của slide.

\section{Tư tưởng cực đại--cực tiểu}

Nhiều bài toán tối ưu có hai mô tả tưởng như đối lập: tìm số đối tượng lớn
nhất có thể chọn, hoặc tìm chướng ngại nhỏ nhất phải loại bỏ. Khi có định lý
cực đại--cực tiểu, hai giá trị này bằng nhau. Ví dụ quen thuộc là luồng cực
đại và cắt nhỏ nhất, ghép cặp lớn nhất và phủ đỉnh nhỏ nhất trong đồ thị hai
phía, hoặc các dạng đường đi ngắn nhất và thế năng đối ngẫu.

\section{Ứng dụng}

Trong bài kiểu \textit{Muddy Fields}, một bài phủ ô có thể chuyển thành ghép
cặp trên đồ thị hai phía. Trong bài có mạng, ta thường dựng cắt để biểu diễn
giới hạn dưới, rồi dùng luồng hoặc đường đi để đạt giới hạn đó. Các slide cũng
nhắc cách dùng thế năng \(d(v)\) trong đường đi ngắn nhất: điều kiện
\[
  d(v)\le d(u)+c(u,v)
\]
có thể xem như một hệ ràng buộc đối ngẫu.

\section{Kết luận}

Khi một bài tối ưu yêu cầu chứng minh rằng đáp án không thể tốt hơn một giá
trị nào đó, hãy tìm một đối tượng đối ngẫu: cắt, phủ, thế năng, màu hoặc hệ
bất đẳng thức. Nếu xây được cả lời giải đạt giá trị ấy, ta có chứng minh tối
ưu gọn và thường có luôn thuật toán.

\end{document}
